\label{sec:implementation}

The \textsc{bisect-apportion} Adams implementation uses the priority:
\[
  P_i^{\text{Ad}}(s) = \frac{p_i}{s}
\]
for $s \geq 1$ (with the first seat guaranteed by the constitutional
floor initialisation).
In integer arithmetic, the comparison is:
$P_i > P_j \iff p_i \cdot s_j > p_j \cdot s_i$,
which uses only 64-bit integers (well within range for 2020 Census data
where $p_i \leq 4 \times 10^7$ and $s_j \leq 52$).

\begin{verbatim}
bisect apportion --year 2020 --method adams
\end{verbatim}

\subsection{European Usage of Adams}

Although Adams has not been used for U.S.\ congressional apportionment,
ceiling rounding (also called the ``method of smallest divisors'')
has been used in some European subnational apportionment contexts.
The Austrian Nationalrat uses a three-stage apportionment process
that includes a ceiling step at the regional level \citep{lijphart1994electoral}.
Italy used ceiling rounding in some historical proportional
representation allocations.
These international uses confirm that Adams is a practically viable
method, not merely a theoretical construct.

\subsection{Test Suite}

\begin{itemize}
  \item \textbf{L0}: $\lceil 1.1 \rceil = 2$, $\lceil 2.0 \rceil = 2$,
    $\lceil 2.1 \rceil = 3$ --- basic ceiling arithmetic.
  \item \textbf{L1}: 3-state example (pop 100, 100, 15; 5 seats):
    Adams with a divisor that produces 5 total seats gives [2, 2, 1]
    (same as HH and Webster for this example).
  \item \textbf{L2}: 2020 Census Adams output matches official HH
    apportionment (since both give the same result for 2020 data)
    at 100\% match rate.
\end{itemize}
